\documentclass[11pt,a4paper]{article}
\usepackage[utf8]{inputenc}
\usepackage[T1]{fontenc}
\usepackage{lmodern}
\usepackage[french]{babel}
\usepackage{amsmath,amssymb,mathtools}
\usepackage{icomma}
\usepackage{xcolor}
\usepackage[most]{tcolorbox}
\usepackage{enumitem}
\usepackage[a4paper,margin=2.2cm,headheight=26pt,headsep=10pt]{geometry}
\usepackage{fancyhdr}
\usepackage[hidelinks]{hyperref}
\definecolor{primary}{RGB}{222,49,99}
\definecolor{primarydark}{RGB}{150,22,55}
\tcbset{boxbase/.style={enhanced,breakable,boxrule=0.9pt,arc=2.5pt,
  left=3mm,right=3mm,top=2.3mm,bottom=2.3mm,fonttitle=\bfseries,coltitle=white}}
\newtcolorbox[auto counter]{exo}[1][]{boxbase,colback=primary!4,colframe=primary,
  title={\textsc{Exercice}~\thetcbcounter\ifx\relax#1\relax\relax\else\textmd{ --- \itshape #1}\fi}}
\newcommand{\R}{\mathbb{R}}
\newcommand{\N}{\mathbb{N}}
\newcommand{\ds}{\displaystyle}
\pagestyle{fancy}\fancyhf{}
\fancyhead[L]{\small\textcolor{primarydark}{\textbf{Thème 1} — Puissances, racines et exposants}}
\fancyhead[R]{\small\textcolor{primarydark}{\textsc{Difficile}}}
\fancyfoot[C]{\small\thepage}
\renewcommand{\headrulewidth}{0.8pt}
\begin{document}

\begin{center}
{\LARGE\bfseries\color{primarydark} Niveau Difficile — Exercices}\\[2pt]
{\color{primary}\rule{0.45\textwidth}{1.2pt}}\\[3pt]
{\small\itshape Niveau concours. Les sous-questions construisent la difficulté pas à pas.}
\end{center}
\vspace{1mm}

\begin{exo}[exposants fractionnaires : simplifier puis évaluer]
Soit $c$ un réel strictement positif et $n\in\N_0$.
\begin{enumerate}[label=\textbf{\alph*)},leftmargin=2.2em,topsep=1pt,itemsep=2pt]
  \item Montre que $\ds\sqrt[n]{\sqrt[n]{c^{\,2n^{2}}}}=c^{2}$, et explique pourquoi le résultat ne dépend pas de $n$.
  \item Simplifie $\ds E=\frac{\sqrt[3]{c}\times\sqrt[6]{c^{5}}}{\sqrt{c}}$ et donne le résultat sous la forme $c^{k}$.
  \item Déduis-en la valeur \emph{exacte} de $E$ lorsque $c=64$ (sans calculatrice).
\end{enumerate}
\end{exo}

\begin{exo}[équations exponentielles : ramener à une base unique]
Résous dans $\R$. À chaque fois, commence par écrire les deux membres avec la même base.
\begin{enumerate}[label=\textbf{\alph*)},leftmargin=2.2em,topsep=1pt,itemsep=2pt]
  \item $2^{2x+1}=2^{-x}$.
  \item $4^{x+1}=8^{x}$.
  \item $\ds\frac{2^{x}\times 2^{x+1}}{2^{3}}=1$.
  \item $\ds\Bigl(\frac{9^{x}}{3}\Bigr)^{2}=3^{\,6x-4}$.
\end{enumerate}
\end{exo}

\begin{exo}[mise en situation : dilutions successives]
Lors d'une série de dilutions, la concentration d'une solution est \emph{divisée par $2$} à chaque
étape. On note $C_0$ la concentration initiale et $C(x)$ la concentration après $x$ étapes
($x\in\N$) : $\ds C(x)=C_0\times\Bigl(\frac12\Bigr)^{x}$.
\begin{enumerate}[label=\textbf{\alph*)},leftmargin=2.2em,topsep=1pt,itemsep=2pt]
  \item Exprime $C(3)$ en fonction de $C_0$, puis simplifie.
  \item Après combien d'étapes la concentration vaut-elle exactement $\ds\frac{C_0}{32}$~? Traduis par une équation exponentielle à même base, puis résous.
  \item À partir de combien d'étapes a-t-on $\ds C(x)\leq\frac{C_0}{64}$~? (Justifie le sens de l'inégalité à l'aide de la monotonie de $x\mapsto\bigl(\tfrac12\bigr)^{x}$.)
\end{enumerate}
\end{exo}

\end{document}
